\section{Calculation of symmetrized and exotic Steinberg maps}

\label{section-4}

\subsection{Conormal direction}

As shown in Theorem \ref{Tp1}, every $K$-orbit in $\Xfv$ takes the form
\[
\Xorbit_\omega=K\cdot([\omega],\flag_0^+,\flag_0^-)
\]
for a matrix $\omega=\begin{pmatrix} \tau_1\\ \tau_2 \end{pmatrix}\in \parameters=\ppermutationsof_{(p,q),r}/\permutationsof{r}$,
where $[\omega]\in\Grass(V,r)$ stands for the image of $\omega$ and $(\flag_0^+,\flag_0^-)\in\Flags(V^+)\times \Flags(V^-)$ is the pair of standard flags.
With the notation of Section \ref{section-1.2} we have
\[\nil([\omega])=\{x\in\gl(V):\mathrm{Im}\,x\subset[\omega]\subset\mathrm{Ker}\,x\},\quad \nil(\flag_0^+)=\fn_p^+,\quad \mbox{and}\quad \nil(\flag_0^-)=\fn_q^+\]
where $\fn_k^+\subset\gl_k(\mC)$ denotes the subalgebra of strictly upper-triangular matrices.
Hence, the conormal bundle to the orbit $\Xorbit_\omega$ is obtained as
\[T^*_{\Xorbit_\omega}\Xfv=K\cdot\{([\omega],\flag_0^+,\flag_0^-,x):x\in\condir_\omega\}\]
where
\begin{equation}
\label{Domega}
\condir_\omega\coloneqq \left\{x=\begin{pmatrix} a & b\\ c & d\end{pmatrix}\in\gl(V):(a,d)\in\fn_p^+\times\fn_q^+,\ \mathrm{Im}\,x\subset[\omega]\subset\mathrm{Ker}\,x\right\}.
\end{equation}
This immediately implies:

\begin{lemma}
\label{L4.1}
Let $\omega\in\parameters$.
Then, $\Phi_\fk(\Xorbit_\omega)$, respectively $\Phi_\fs(\Xorbit_\omega)$,
is characterized as being the unique $K$-orbit of $\nilpotentsof{\fk}$, resp. $\nilpotentsof{\fs}$, which intersects the space $\{x_\fk:x\in\condir_\omega\}$, resp. $\{x_\fs:x\in\condir_\omega\}$, along a dense open subset.
\end{lemma}

For the computation of the maps $\Phi_\fk$ and $\Phi_\fs$, we need a more detailed description of the conormal direction $\condir_\omega$.
We use the following notation: if $a$ is a $k\times \ell$ matrix, then for every subsets $R\subset\{1,\ldots,k\}$ and $S\subset\{1,\ldots,\ell\}$, we denote by $\subm{a}{R,S}$ the submatrix of $a$ formed by the coefficients $a_{i,j}$ with $i\in R$, $j\in S$, and we view it as a linear map from $\mC^S$ to $\mC^R$. Recall from Section \ref{section-2.1}
that $\omega$ gives rise to decompositions $\intp\coloneqq \{1,\ldots,p\}=I\cup L\cup L'$ and $\intq\coloneqq \{1,\ldots,q\}=J\cup M\cup M'$ and to a bijection $\sigma:J\to I$ which we view as a linear map from $\mC^J$ to $\mC^I$.

\begin{lemma}
\label{L4.2}
A matrix $x=\begin{pmatrix} a & b\\ c & d \end{pmatrix}$ \textup{(}with $a\in\fn_p^+$, $d\in\fn_q^+$\textup{)} belongs to $\mathcal{D}_\omega$ if and only if it satisfies the following conditions:
\begin{eqnarray}
&& \left\{\begin{array}{c} \subm{a}{\intp,L}=0,\ \subm{c}{\intq,L}=0,\ \subm{b}{\intp,M}=0,\ \subm{d}{\intq,M}=0,\\[1.5mm]
{} \subm{b}{\intp,J}=-\subm{a}{\intp,I}\sigma,\ \subm{d}{\intq,J}=-\subm{c}{\intq,I}\sigma, \end{array}\right. \label{4.1}\\[2mm]
&& \left\{\begin{array}{c} \subm{a}{L',\intp}=0,\ \subm{b}{L',\intq}=0,\ \subm{c}{M',\intp}=0,\ \subm{d}{M',\intq}=0,\\[1.5mm]
{} \subm{a}{I,\intp}=\sigma \subm{c}{J,\intp},\ \subm{b}{I,\intq}=\sigma\subm{d}{J,\intq}.
\end{array}\right. \label{4.2}
\end{eqnarray}
\end{lemma}

\begin{proof}
Let $(e_1^+,\ldots,e_p^+)$ be the standard basis of $V^+=\mC^p\times\{0\}^q$ and let $(e_1^-,\ldots,e_q^-)$ be the standard basis of $V^-=\{0\}^p\times\mC^q$. Then
\begin{equation}
\label{4.3}
[\omega]=\Span{e_i^+:i\in L;\quad e_j^-:j\in M;\quad e_{\sigma(j)}^++e_j^-:j\in J}.
\end{equation}
For every matrix $x$ such that $a\in\fn_p^+$ and $d\in\fn_q^+$, we have $x\in\condir_\omega$ if and only if $\mathrm{Im}\,x\subset[\omega]\subset \mathrm{Ker}\,x$, and in view of \eqref{4.3} the second inclusion is equivalent to \eqref{4.1} while the first inclusion is equivalent to \eqref{4.2}.
\end{proof}

Figure \ref{figure4} illustrates the form of the elements in the conormal direction $\mathcal{D}_\omega$. The matrix is represented blockwise with blocks indicating the submatrices $\subm{X}{R,S}$ relative to the subsets
$R,S\in\{I,J,L,L',M,M'\}$. Note that in the decompositions $I\cup L\cup L'=\intp$ and $J\cup M\cup M'=\intq$ the subsets are not consecutive,
and the matrix is thus represented modulo permutation within the rows and the columns. In particular, it is required in addition that the blocks $a$ and $d$ be strictly upper triangular.
\begin{figure}[h]
\[
\begin{blockarray}{ccccccc}
{} & {\quad I\quad} & {\quad L\quad} & {\quad L'\quad} & {\quad J\quad} & {\quad M\quad} & {\quad M'\quad}\\
\begin{block}{c(c|c|c||c|c|c)}
I & \sigma c_1 & 0 & \sigma c_2 & -\sigma c_1\sigma & 0 & \sigma d_2\\
\cline{2-7}
L & a_3 & 0 & a_4 & -a_3\sigma & 0 & b_4\\
\cline{2-7}
L' & 0 & 0 & 0 & 0 & 0 & 0\\
\cline{2-7}\\[-12pt] \cline{2-7}
J & c_1 & 0 & c_2 & -c_1\sigma & 0 & d_2\\
\cline{2-7}
M & c_3 & 0 & c_4 & -c_3\sigma & 0 & d_4\\
\cline{2-7}
M' & 0 & 0 & 0 & 0 & 0 & 0\\
\end{block}
\end{blockarray}
\]
\caption{Form of the elements $x=\left({}^{a \ b}_{c \ d}\right)$ (with $a\in\fn_p^+$, $d\in\fn_q^+$) belonging to the conormal direction $\condir_\omega$.}
\label{figure4}
\end{figure}

\subsection{A review of the orbit sets, and an involution}
As explained in Section \ref{section-2.3}, the orbits of $K$ in the nilpotent cone
\[
\nilpotentsof{\fk}\subset\fk=\left\{\begin{pmatrix} a & 0\\ 0 & d \end{pmatrix}: (a,d)\in \gl_p(\mC)\times\gl_q(\mC)\right\}
\]
are parametrized by pairs of partitions $(\lambda,\mu)\in\partitionsof{p}\times\partitionsof{q}$, and we denote by $\Nkorbit_{\lambda,\mu}$ the orbit corresponding to the pair $(\lambda,\mu)$.
Note that
\begin{equation}
\label{4.5new}
\begin{pmatrix} a & 0\\ 0 & d\end{pmatrix}\in\Nkorbit_{\lambda,\mu}\Longleftrightarrow \begin{pmatrix} d & 0\\ 0 & a \end{pmatrix}\in\Nkorbit_{\mu,\lambda},
\end{equation}
though here the notation $\Nkorbit_{\mu,\lambda}$ refers to an orbit of $K^*\coloneqq \GL_q(\mC)\times\GL_p(\mC)$ on $\gl_q(\mC)\times\gl_p(\mC)$.

Recall also from Section \ref{section-2.3} that the orbits of $K$ in the nilpotent cone
\[
\nilpotentsof{\fs}\subset\fs=\left\{\begin{pmatrix} 0 & b\\ c & 0 \end{pmatrix}: (b,c)\in \Mat_{p,q}(\mC)\times\Mat_{q,p}(\mC)\right\}
\]
are parametrized by signed Young diagrams of signature $(p,q)$, and we denote by $\Nsorbit_\Lambda$ the orbit corresponding to $\Lambda$. We have
\begin{equation}
\label{4.6}
\begin{pmatrix} 0 & b\\ c & 0 \end{pmatrix}\in\Nsorbit_\Lambda\Longleftrightarrow \begin{pmatrix} 0 & c\\ b & 0 \end{pmatrix}\in\Nsorbit_{\Lambda^*}
\end{equation}
where $\Lambda^*$ denotes the signed Young diagram of signature $(q,p)$ obtained from $\Lambda$ by switching the $+$'s and the $-$'s, and $\Nsorbit_{\Lambda^*}$ is an orbit of the group $K^*$.

As shown in Theorem \ref{Tp1}, the orbits of $K$ in
\[
\Xfv=\Grass(V,r)\times\Flags(V^+)\times\Flags(V^-)
\]
are parametrized by the elements of $\parameters$, and
$\Xorbit_\omega$ is the orbit corresponding to $\omega$. If
$\omega=\begin{pmatrix} \tau_1\\ \tau_2 \end{pmatrix}$ is an
element of
$\parameters=\ppermutationsof_{(p,q),r}/\permutationsof{r}$, then
$\omega^*\coloneqq \begin{pmatrix} \tau_2\\ \tau_1 \end{pmatrix}$ is an
element of
$\parameters^*\coloneqq \ppermutationsof_{(q,p),r}/\permutationsof{r}$ which
thus yields an orbit $\Xorbit_{\omega^*}$ of $K^*$ in a suitable
multiple flag variety $\Xfv^*$. The graphic representation
$\mathcal{G}(\omega^*)$ of $\omega^*$ is obtained from
$\mathcal{G}(\omega)$ by switching the two rows of vertices, i.e.
by relabeling every vertex $i^+$ (resp. $j^-$) as $i^-$ (resp.
$j^+$). This implies that, if $(I,J,L,L',M,M',\sigma)$ are the data
corresponding to $\omega$ in the sense of Section \ref{section-2.1},
then the relevant data for $\omega^*$ are
\begin{equation}
\label{4.7}
(I^*,J^*,L^*,L'^*,M^*,M'^*,\sigma)=(J,I,M,M',L,L',\sigma^{-1}).
\end{equation}

Finally, by the description of the conormal direction in \eqref{Domega}, we have
\begin{equation}
\label{4.8}
\begin{pmatrix} a & b\\ c & d \end{pmatrix}\in\condir_\omega\Longleftrightarrow
\begin{pmatrix} d & c\\ b & a \end{pmatrix}\in\condir_{\omega^*}.
\end{equation}

The observations made in \eqref{4.5new}, \eqref{4.6}, \eqref{4.8} combined with Lemma \ref{L4.1} yield the following statement:

\begin{lemma}
\label{L4.3}
\begin{enumerate}
\item If $\Phi_\fk(\Xorbit_\omega)=\Nkorbit_{\lambda,\mu}$, then $\Phi_\fk(\Xorbit_{\omega^*})=\Nkorbit_{\mu,\lambda}$.
\item If $\Phi_\fs(\Xorbit_\omega)=\Nsorbit_\Lambda$, then $\Phi_\fs(\Xorbit_{\omega^*})=\Nsorbit_{\Lambda^*}$.
\end{enumerate}
\end{lemma}

There is an abuse of notation in that statement, since we use the notation $\Phi_\fk$ and $\Phi_\fs$ to designate also the symmetrized and exotic Steinberg maps relative to the symmetric pair $(G,K^*)=(\GL_{p+q}(\mC),\GL_q(\mC)\times\GL_p(\mC))$.

\subsection{Symmetrized Steinberg map $\Phi_\fk$}

For a permutation $w\in\permutationsof{k}$ we consider the space
\[
\fn_k^+\cap({}^w\fn_k^+)\coloneqq \fn_k^+\cap(\mathrm{Ad}(w)(\fn_k^+))=\{a\in\fn_k^+:w^{-1}aw\in\fn_k^+\}.
\]
The following is a well-known fact from classical Steinberg theory.

\begin{theorem}[\cite{Steinberg-occurrence}]
\label{T-Steinberg}
The unique nilpotent orbit $\mathcal{O}_\lambda\subset \gl_k(\mC)$ which intersects the space $\fn_k^+\cap({}^w\fn_k^+)$ along a dense open subset is the one corresponding to the Young diagram
$\lambda=\shape(\RS_1(w))=\shape(\RS_2(w))$, where $(\RS_1(w),\RS_2(w))$ denotes the pair of Young tableaux associated to $w$ via the Robinson--Schensted correspondence.
\end{theorem}

In our situation, we consider the permutations $w_{\fk,+}\in\permutationsof{p}$ and $w_{\fk,-}\in\permutationsof{q}$
of \eqref{wk+} and \eqref{wk-}, associated to an element $\omega\in\parameters$.

\begin{lemma}
\label{L4.5}
%\begin{enumerate}
%\item
Let $a\in\fn_p^+$. The following conditions are equivalent:
\begin{enumerate}[label=(\arabic*)]
\item \label{it1vvt1}There are matrices $b,c,d$ such that $x\coloneqq \begin{pmatrix} a & b\\ c & d\end{pmatrix}\in\condir_\omega$;
\item \label{it1vvt2} $\sigma^{-1}\subm{a}{I,I}\sigma$ is strictly upper triangular, $\subm{a}{L',\intp}=0$, and $\subm{a}{\intp,L}=0$;
\item\label{it1vvt3} $a\in \fn_p^+\cap({}^{w_{\fk,+}}\fn_p^+)$.
\end{enumerate}
\end{lemma}

\begin{proof}
The equivalence between conditions~\ref{it1vvt1} and ~\ref{it1vvt2} is implied by Lemma \ref{L4.2} (see Figure \ref{figure4}).
Given $a\in\fn_p^+$, condition~\ref{it1vvt2}  is equivalent to:
\begin{equation*}
\left(1\leq i<j\leq p\quad\mbox{and}\quad\left\{\begin{array}{l}
i\in L'\\ \mbox{or}\quad j\in L\\
\mbox{or}\quad(i,j\in I\ \ \mbox{and}\ \ \sigma^{-1}(i)>\sigma^{-1}(j))
\end{array}\right.\right)
\quad\implies\quad a_{i,j}=0.
\end{equation*}
Condition~\ref{it1vvt3} is equivalent to:
\begin{equation*}
1\leq i<j\leq p\quad \mbox{and}\quad w_{\fk,+}^{-1}(i)> w_{\fk,+}^{-1}(j)\quad \implies\quad  a_{i,j}=0.
\end{equation*}
By definition of $w_{\fk,+}$ (see \eqref{wk+}), for $1\leq i<j\leq p$, we have
\[w_{\fk,+}^{-1}(i)> w_{\fk,+}^{-1}(j)\quad\iff\quad \left\{\begin{array}{l}
i\in L'\\ \mbox{or}\quad j\in L\\
\mbox{or}\quad(i,j\in I\ \ \mbox{and}\ \ \sigma^{-1}(i)>\sigma^{-1}(j)).
\end{array}\right.\]
Therefore, conditions~\ref{it1vvt2} and~\ref{it1vvt3} are equivalent.
\end{proof}


\begin{proof}[Proof of Theorem \ref{T2}\ref{it2vvt1}]
Let
\[\Phi_\fk(\Xorbit_\omega)=\Nkorbit_{\lambda,\mu}=\left\{\begin{pmatrix} a & 0\\ 0 & d \end{pmatrix}:(a,d)\in\mathcal{O}_\lambda\times\mathcal{O}_\mu\subset\gl_p(\mC)\times\gl_q(\mC)\right\}.\]
By Lemmas \ref{L4.1} and \ref{L4.5}, the nilpotent orbit $\mathcal{O}_\lambda\subset\gl_p(\mC)$ is characterized as being the unique $\GL_p(\mC)$-orbit which intersects the space
\[\left\{a\in\fn_p^+:\exists b,c,d\mbox{ such that }\begin{pmatrix} a & b\\ c & d \end{pmatrix}\in\condir_\omega\right\}=\fn_p^+\cap({}^{w_{\fk,+}}\fn_p^+)\]
along a dense open subset. By Theorem \ref{T-Steinberg}, this implies that
\[\lambda=\shape(\RS_1(w_{\fk,+})).\]
By \eqref{4.7} and Lemma \ref{L4.3}, we also deduce that
\[\mu=\shape(\RS_1(w_{\fk,-})).\]
The proof of Theorem \ref{T2}\ref{it2vvt1} is complete.
\end{proof}

\subsection{Exotic Steinberg map $\Phi_\fs$}

We introduce notation which extend our notation on matrices.
Given subsets $R,S$ of integers, we let $\Mat_{R,S}(\mC)$ denote the space of linear homomorphisms $x:\mC^S\to \mC^R$, which can be viewed as well as matrices of coefficients $(x_{i,j})_{(i,j)\in R\times S}$.
Let $\fn_R^+\subset\Mat_{R,R}(\mC)$ be the subspace of endomorphisms that are strictly upper triangular as matrices. %in the standard basis of $\mC^R$.

If $R,S$ are respectively subsets of $R',S'$, we denote by
\begin{equation}
\label{eta}
\eta_{R,S}^{R',S'}:\Mat_{R,S}(\mC)\to\Mat_{R',S'}(\mC),\quad x\mapsto \hat{x}=(\hat{x}_{i,j})_{(i,j)\in R'\times S'}
\end{equation}
the linear morphism which maps a matrix $x$ to its {\em extension by zero} given by
$\hat{x}_{i,j}=x_{i,j}$ if $(i,j)\in R\times S$ and $\hat{x}_{i,j}=0$ otherwise.

A bijection $w:S\to R$ yields an element of $\Mat_{R,S}(\mC)$ also denoted by $w$ by abuse of notation. In addition, through the Robinson--Schensted algorithm, $w$ gives rise to a pair of Young tableaux $(\RS_1(w),\RS_2(w))$ of the same shape, whose respective sets of entries are $R$ and $S$.

The following is a reformulation of Theorem \ref{T-Steinberg}.

\begin{prop}
\label{P4.6}
Given a bijection $w: S\to R$, the Jordan normal form of a general element $x$ in the space
\[
\fn_R^+\cap({}^w\fn_S^+)\coloneqq \{x\in\fn_R^+:w^{-1}xw\in\fn_S^+\}
\]
is given by the Young diagram $\shape(\RS_1(w))=\shape(\RS_2(w))$. In other words, for all $k\geq 1$, $\dim \ker x^k$ is the number of boxes in the first $k$ columns of $\RS_1(w)$.\end{prop}

Take an element $\omega\in\parameters$ with corresponding data $(I,J,L,L',M,M',\sigma)$. As in Section \ref{section-2.3}, we write $J=\{j_1<\ldots<j_k\}$, $L'=\{\ell'_1<\ldots<\ell'_{s'}\}$, $M=\{m_1<\ldots<m_t\}$.
Moreover, we denote
\[
S=J\cup M\cup \{q+1,\ldots,q+s'\}\quad\mbox{and}\quad
 R=\{-t,\ldots,-1\}\cup I\cup L'.
\]
We will consider the bijection $w\coloneqq w_{\fs,+}:S\to R$ defined in \eqref{ws+}.


We denote $\varsigma=\eta_{I,J}^{I\cup L',J\cup M}(\sigma)$ and $\tau=\eta_{I,J}^{\intp,\intq}(\sigma)=\eta_{I\cup L',J\cup M}^{\intp,\intq}(\varsigma)$ (see \eqref{eta}), that is,
\begin{equation}
\label{zeta-tau}
\varsigma=\begin{blockarray}{ccc}
{} & {\ J\ } & {\ M\ }\\
\begin{block}{c(c|c)}
I & \sigma & 0\\
\cline{2-3}
L' & 0 & 0\\
\end{block}
\end{blockarray}\,,\qquad \tau=\begin{blockarray}{cccc}
{} & {\ J\ } & {\ M\ } & {\ M'\ }\\
\begin{block}{c(c|c|c)}
I & \sigma & 0 & 0\\
\cline{2-4}
L & 0 & 0 & 0\\
\cline{2-4}
L' & 0 & 0 & 0\\
\end{block}
\end{blockarray}
\end{equation}
(after changing the order of rows and columns).
%Note also that $w$ above is an element of $\Mat_{R,S}(\mC)$.

\begin{lemma}
\label{L4.8}
Let $x=\begin{pmatrix} a & b\\ c & d \end{pmatrix}\in\condir_\omega$, so that $x_\fk=\begin{pmatrix} a & 0\\ 0 & d\end{pmatrix}$ and $x_\fs=\begin{pmatrix} 0 & b\\ c & 0\end{pmatrix}$. Then, for every $m\geq 0$, we have
\[
(x_\fs)^{2m}=(-1)^m(x_\fk)^{2m}=(-1)^m\begin{pmatrix} a^{2m} & 0\\ 0 & d^{2m} \end{pmatrix}
\ \ \mbox{and}\ \
(x_\fs)^{2m+1}=(-1)^m\begin{pmatrix} 0 & *\\ (c\tau)^{2m}c & 0 \end{pmatrix}
\]
\textup{(}where the symbol $*$ stands for some matrix in $\Mat_{p,q}(\mC)$, whose precise description is not needed\textup{)}.
\end{lemma}

\begin{proof}
Every element $x\in\condir_\omega$ is such that $x^2=0$. On the other hand, we can see that $(x^2)_\fk=(x_\fk)^2+(x_\fs)^2$. The formula for $(x_\fs)^{2m}$ ensues.

It readily follows that
\[(x_\fs)^{2m+1}=(-1)^m\begin{pmatrix} 0 & a^{2m}b\\ d^{2m}c & 0 \end{pmatrix}.\]
For every $\ell\geq 1$, using the notation of Figure \ref{figure4},
we can see that $(-1)^\ell d^\ell c$ is the matrix (written
blockwise)
\begin{equation}
\label{powers1}
\begin{blockarray}{cccc}
{} & {\quad I\quad} & {\quad L\quad} & {\quad L'\quad}\\
\begin{block}{c(c|c|c)}
J & (c_1\sigma)^\ell c_1 & 0 & (c_1\sigma)^\ell c_2\\
\cline{2-4}
M & c_3\sigma(c_1\sigma)^{\ell-1}c_1 & 0 & c_3\sigma(c_1\sigma)^{\ell-1}c_2\\
\cline{2-4}
M' & 0 & 0 & 0\\
\end{block}
\end{blockarray}
\,,
\end{equation}
and this coincides with $(c\tau)^\ell c$.
Letting $\ell=2m$, this yields the formula claimed for $(x_\fs)^{2m+1}$ in the case of $m\geq 1$. The claimed formula is immediate if $m=0$.
\end{proof}

We consider the following extension-by-zero mappings:
\[
\xymatrix{
& \ar_{\eta_1\coloneqq \eta_{J\cup M,I\cup L'}^{\intq,\intp}\quad}[ld] \Mat_{J\cup M,I\cup L'}(\mC) \ar^{\quad\eta_2\coloneqq \eta_{J\cup M,I\cup L'\\ }^{S,R}}[rd] &\\
\Mat_{\intq,\intp}(\mC)=\Mat_{q,p}(\mC) &  & \Mat_{S,R}(\mC).}
\]
In addition, we consider the subspace $\Gamma\coloneqq \{\gamma\in \Mat_{J\cup M,I\cup L'}(\mC):\varsigma \gamma\in\fn_{I\cup L'}^+,\ \gamma\varsigma\in\fn_{J\cup M}^+\}$ (where $\varsigma$ is as in \eqref{zeta-tau}).

\begin{lemma}
\label{L4.9}
The maps $\eta_1$ and $\eta_2$ restrict to linear isomorphisms
\[(\eta_1)|_{\Gamma}:\Gamma\stackrel{\sim}{\longrightarrow} \left\{ c\in \Mat_{q,p}(\mC):\exists a,b,d\ \mbox{such that}\ \begin{pmatrix} a & b\\ c & d \end{pmatrix}\in\condir_\omega \right\}\]
and\[(\eta_2)|_\Gamma:\Gamma\stackrel{\sim}{\to} \{z\in\Mat_{S,R}(\mC):wz\in\fn_R^+,\ zw\in\fn_S^+\}=\{z\in\Mat_{S,R}(\mC):wz\in\fn_R^+\cap({}^w\fn_S^+)\},\]
where $w=w_{\fs,+}$ is the bijection defined in \eqref{ws+}.
Moreover, for every $\gamma\in\Gamma$ such that $c=\eta_1(\gamma)$ and $z=\eta_2(\gamma)$, we have
\[\eta_1((\gamma\varsigma)^{\ell}\gamma)=(c\tau)^\ell c\quad\mbox{and}\quad \eta_2((\gamma\varsigma)^\ell\gamma)=(zw)^\ell z=w^{-1}(wz)^{\ell+1}\quad\mbox{for all $\ell\geq 0$},\]
where $\varsigma$ and $\tau$ are given by \eqref{zeta-tau}.
\end{lemma}

\begin{proof}
The space $\Gamma$ consists of the matrices
\begin{equation}
\label{gamma}
\gamma=\begin{blockarray}{ccc}
{} & {\quad I\quad} & {\quad L'\quad}\\
\begin{block}{c(c|c)}
J & c_1 & c_2\\
\cline{2-3}
M & c_3 & c_4\\ \end{block}
\end{blockarray}
\,\in \Mat_{J\cup M,I\cup L'}(\mC)
\end{equation}
(written blockwise) such that
\begin{equation}
\label{zeta-gamma}
\varsigma \gamma=\begin{blockarray}{ccc}
{} & {\quad I\quad} & {\quad L'\quad}\\
\begin{block}{c(c|c)}
I & \sigma c_1 & \sigma c_2\\
\cline{2-3}
L' & 0 & 0\\ \end{block}
\end{blockarray}
\qquad\mbox{and}\qquad \gamma\varsigma=\begin{blockarray}{ccc}
{} & {\quad J\quad} & {\quad M\quad}\\
\begin{block}{c(c|c)}
J & c_1\sigma & 0\\
\cline{2-3}
M & c_3\sigma & 0\\ \end{block}
\end{blockarray}\end{equation}
are strictly upper triangular (before rearranging the rows and the columns).
Moreover, for every $\ell\geq 1$, we have
\begin{equation}
\label{powers2}
(\gamma\varsigma)^\ell\gamma=\begin{blockarray}{ccc}
{} & {\quad I\quad} & {\quad L'\quad}\\
\begin{block}{c(c|c)}
J & (c_1\sigma)^\ell c_1 & (c_1\sigma)^\ell c_2\\
\cline{2-3}
M & c_3\sigma(c_1\sigma)^{\ell-1}c_1 & c_3\sigma(c_1\sigma)^{\ell-1}c_2\\ \end{block}
\end{blockarray}\,.
\end{equation}

One can see from Figure \ref{figure4} that the space
\[
\mathcal{C}\coloneqq \left\{ c\in\Mat_{q,p}(\mC):\exists a,b,d\ \mbox{such that}\ \begin{pmatrix} a & b\\ c & d \end{pmatrix}\in\condir_\omega\right\}
\]
consists of the matrices of the form
\[
c=\begin{blockarray}{cccc}
{} & {\quad I\quad} & {\quad L\quad} & {\quad L'\quad}\\
\begin{block}{c(c|c|c)}
J & c_1 & 0 & c_2\\
\cline{2-4}
M & c_3 & 0 & c_4\\
\cline{2-4}
M' & 0 & 0 & 0\\
\end{block}
\end{blockarray}
\]
such that the matrices
\begin{eqnarray*}
 & \begin{blockarray}{cccc}
{} & {\quad I\quad} & {\quad L\quad} & {\quad L'\quad}\\
\begin{block}{c(c|c|c)}
I & \sigma c_1 & 0 & \sigma c_2\\
\cline{2-4}
L & 0 & 0 & 0\\
\cline{2-4}
L' & 0 & 0 & 0\\
\end{block}
\end{blockarray}\in\Mat_p(\mC)\\
\mbox{and}&
\begin{blockarray}{cccc}
{} & {\quad J\quad} & {\quad M\quad} & {\quad M'\quad}\\
\begin{block}{c(c|c|c)}
J & -c_1\sigma & 0 & 0\\
\cline{2-4}
M & -c_3\sigma & 0 & 0\\
\cline{2-4}
M' & 0 & 0 & 0\\
\end{block}
\end{blockarray}\in\Mat_q(\mC)
\end{eqnarray*}
are strictly upper triangular (before rearranging the rows and the columns). From the above description of $\Gamma$, these conditions are equivalent to having that $c=\eta_1(\gamma)$ for some $\gamma\in\Gamma$. Hence $\mathcal{C}=\eta_1(\Gamma)$, and this implies that $\eta_1$ restricts to a linear isomorphism $(\eta_1)|_\Gamma:\Gamma\to \mathcal{C}$.
Moreover, by comparing the expressions of $(c\tau)^\ell c$ and $(\gamma\varsigma)^\ell\gamma$ given in \eqref{powers1} and \eqref{powers2}, respectively, one can see that the equality $\eta_1((\gamma\varsigma)^\ell\gamma)=(c\tau)^\ell c$ holds for all $\ell\geq 0$ whenever $\eta_1(\gamma)=c$.

We have $R=R'\cup I\cup L'$ and $S=J\cup M\cup S'$ where $R'\coloneqq \{-i\}_{i=1}^t$
and $S'\coloneqq \{q+j\}_{j=1}^{s'}$.
Let $z$ be an element in $\Mat_{S,R}(\mC)$, and let us write it (blockwise) as
\begin{equation}
\label{z}
z=\begin{blockarray}{cccc}
{} & {\quad R'\quad} & {\quad I\quad} & {\quad L'\quad}\\
\begin{block}{c(c|c|c)}
J & z_1 & z_2 & z_3\\
\cline{2-4}
M & z_4 & z_5 & z_6\\
\cline{2-4}
S' & z_7 & z_8 & z_9\\
\end{block}
\end{blockarray}\,.\end{equation}
The bijection $w=w_{\fs,+}:S\to R$ of \eqref{ws+}) can be written in the following matrix form
\[
w=\begin{blockarray}{cccc}
{} & {\quad J\quad} & {\quad M\quad} & {\quad S'\quad}\\
\begin{block}{c(c|c|c)}
R' & 0 & \sigma_0 & 0\\
\cline{2-4}
I & \sigma & 0 & 0\\
\cline{2-4}
L' & 0 & 0 & \sigma_0'\\
\end{block}
\end{blockarray}
\]
where the block $\sigma_0\in\Mat_{R',M}(\mC)$ is yielded by the unique decreasing bijection $M\to R'$; this corresponds to a block with $1$'s on the antidiagonal and $0$'s elsewhere. The block $\sigma_0'\in\Mat_{L',S'}(\mC)$ is defined in the same way.
We get
\[ wz=\begin{blockarray}{cccc}
{} & {\quad R'\quad} & {\quad I\quad} & {\quad L'\quad}\\
\begin{block}{c(c|c|c)}
R' & \sigma_0 z_4 & \sigma_0 z_5 & \sigma_0 z_6\\
\cline{2-4}
I & \sigma z_1 & \sigma z_2 & \sigma z_3\\
\cline{2-4}
L' & \sigma_0' z_7 & \sigma_0' z_8 & \sigma_0' z_9\\
\end{block}
\end{blockarray} \quad
\mbox{and}\quad
zw=\begin{blockarray}{cccc}
{} & {\quad J\quad} & {\quad M\quad} & {\quad S'\quad}\\
\begin{block}{c(c|c|c)}
J & z_2\sigma & z_1\sigma_0 & z_3\sigma_0'\\
\cline{2-4}
M & z_5\sigma & z_4\sigma_0 & z_6\sigma_0'\\
\cline{2-4}
S' & z_8\sigma & z_7\sigma_0 & z_9\sigma_0'\\
\end{block}
\end{blockarray}\,.
\]
Note that we have $i<j$ whenever $(i,j)\in R'\times(I\cup L')$. We have also $i<j$ whenever $(i,j)\in(J\cup M)\times S'$. We obtain that $(wz,zw)\in\fn_R^+\times \fn_S^+$ if and only if the following conditions are satisfied:
\[
\left\{\begin{array}{ll}
\mbox{$\bullet$ the blocks $z_1,z_4,z_7,z_8,z_9$ are zero},\\
\parbox{11cm}{$\bullet$ the submatrices $\begin{blockarray}{ccc}
{} & {\quad I\quad} & {\quad L'\quad}\\
\begin{block}{c(c|c)}
I & \sigma z_2 & \sigma z_3\\
\cline{2-3}
L' & 0 & 0\\
\end{block}
\end{blockarray}$
\ and $\begin{blockarray}{ccc}
{} & {\quad J\quad} & {\quad M\quad}\\
\begin{block}{c(c|c)}
J & z_2\sigma & 0\\
\cline{2-3}
M & z_5\sigma & 0\\
\end{block}
\end{blockarray}$
\ are strictly upper triangular.}
\end{array}\right.
\]
Combining these observations with the description of the space $\Gamma$ given above, we conclude that the elements in the space $\mathcal{Z}\coloneqq \{z\in\Mat_{S,R}(\mC):wz\in\fn_R^+,\ zw\in\fn_S^+\}$ are exactly of the form $z=\eta_2(\gamma)$ for $\gamma\in\Gamma$ (see \eqref{gamma}--\eqref{zeta-gamma}) and, therefore, $\eta_2$ restricts to an isomorphism $(\eta_2)|_\Gamma:\Gamma\to\mathcal{Z}$ as asserted.

Finally, let $z=\eta_2(\gamma)$ for $\gamma\in\Gamma$ written as in \eqref{gamma}.
In the notation of \eqref{z}, this means that $z_2=c_1$, $z_3=c_2$, $z_5=c_3$, $z_6=c_4$, and the other blocks of $z$ are zero. In view of the expression for $zw$ and $(\gamma\varsigma)^\ell\gamma$ given in \eqref{powers2}, this yields
\[(zw)^\ell z=\begin{blockarray}{cccc}
{} & {\quad R'\quad} & {\quad I\quad} & {\quad L'\quad}\\
\begin{block}{c(c|c|c)}
J & 0 & (c_1\sigma)^\ell c_1 & (c_1\sigma)^{\ell}c_2\\
\cline{2-4}
M & 0 & c_3\sigma(c_1\sigma)^{\ell-1}c_1 & c_3\sigma(c_1\sigma)^{\ell-1}c_2\\
\cline{2-4}
S' & 0 & 0 & 0\\
\end{block}
\end{blockarray}\,=\eta_2((\gamma\varsigma)^\ell\gamma)\]
for all $\ell\geq 1$. The lemma is proved.
\end{proof}

\begin{proof}[Proof of Theorem \ref{T2}\ref{it2vvt2}]
Let $\Phi_\fk(\Xorbit_\omega)=\Nkorbit_{\lambda,\mu}$ and $\Phi_\fs(\Xorbit_\omega)=\Nsorbit_\Lambda$.
Let
\[x=\begin{pmatrix}a & b\\ c & d\end{pmatrix}\]
be a general element of $\mathcal{D}_\omega$, so that $x_\fk\in\Nkorbit_{\lambda,\mu}$ and $x_\fs\in\Nsorbit_\Lambda$
(see Lemma \ref{L4.1}).
It follows from Lemma \ref{L4.8} that, for every $\ell\geq 1$,
the number $\nrplus{\Lambda}{\ell}$ of $+$'s in the first $\ell$ columns of the signed Young diagram $\Lambda$ is equal to
\[\dim\ker a^{2m}\mbox{ if $\ell=2m$ is even},\quad\mbox{respectively}\quad \dim \ker(c\tau)^{2m} c\mbox{ if $\ell=2m+1$ is odd}\]
(with $\tau$ from \eqref{zeta-tau}).
If $\ell=2m$ is even, Theorem \ref{T2}\,(1) shows that this number is equal to the number $\nrboxes{\lambda}{2m}$ of boxes in the first $2m$ columns of $\lambda$.
This confirms the first assertion made in Theorem \ref{T2}\ref{it2vvt2}\,(a).

Now assume that $\ell=2m+1$ is odd. By Lemma \ref{L4.9}, we have $c=\eta_1(\gamma)$ with $\gamma\in\Gamma$. Moreover, we can assume that $z\coloneqq \eta_2(\gamma)$ is general in $\{z\in\Mat_{S,R}(\mC):wz\in\fn_R^+\cap({}^w\fn_S^+)\}$,
with $w=w_{\fs,+}$ as in \eqref{ws+}. Hence, by Proposition \ref{P4.6}, this element $z$ is such that
\[
\dim \ker (wz)^{2m+1}=\nrboxes{\lambda'}{2m+1}
\]
with $\lambda'=\shape(\RS_1(w))$.
Note also that $\dim\ker(wz)^{2m+1}=\dim\ker(zw)^{2m}z$. In view of the previous observations, and by Lemma \ref{L4.9}, we get
\begin{eqnarray*}
\nrplus{\Lambda}{2m+1} & = & \dim \ker(c\tau)^{2m} c\\ & = & \card{L}+\dim \ker (\gamma\varsigma)^{2m}\gamma\\
 & = & \card{L}+\dim \ker(zw)^{2m}z-\card{R'}\\
 & = & s-t+\nrboxes{\lambda'}{2m+1}.
\end{eqnarray*}
This establishes the claim in Theorem \ref{T2}\ref{it2vvt2}\,(b) regarding the number of $+$'s. The formulas regarding the number of $-$'s stated in Theorem \ref{T2}\ref{it2vvt2}\,(a)--(b) can now be deduced, by invoking Lemma \ref{L4.3} and taking \eqref{4.7} into account. The proof of Theorem \ref{T2}\ref{it2vvt2} is complete.
\end{proof}

